\section{Layout du circuit imprimé}
\label{sec:layout}

Le schéma électrique fixe les connexions à établir, le layout détermine la manière dont elles se comportent une fois réalisées. Une piste n'est pas un fil idéal, elle possède une inductance et une capacité, elle rayonne et elle capte, et la disposition des composants conditionne la longueur des boucles de courant et la qualité des références de tension. Cette section décrit les décisions de placement et de routage qui découlent de ces contraintes, puis la démarche de vérification appliquée avant la commande des cartes. Les vues complètes des couches, de la sérigraphie et du plan de perçage figurent en annexe~\ref{ann:layout}.

\subsection{Stack-up et contraintes}

Une partie du cadre n'a pas fait l'objet d'un choix. Le cahier des charges impose un circuit imprimé quatre couches, et le projet Altium, le jeu de règles de conception et le stack-up ont été fournis par l'\gls{heig-vd}, calés sur les tolérances de la classe 6C d'Eurocircuits~\cite{eurocircuits_classes}. Le stack-up retenu place le signal sur le Top Layer, un plan de masse en Int1, un plan de puissance en Int2 et le signal sur le Bottom Layer, pour une épaisseur totale de 1,629~mm. Le diélectrique séparant le Top Layer du plan de masse mesure 0,360~mm, et le noyau séparant les deux plans internes mesure 0,710~mm.

Ces deux distances pèsent sur le routage. Une piste du Top Layer se trouve relativement loin de son plan de référence, ce qui élargit la boucle formée par le signal et son courant de retour et impose, pour une impédance donnée, des pistes plus larges que sur un stack-up où le plan de masse est immédiatement sous la couche de signal. L'épaisseur du noyau signifie par ailleurs que la capacité répartie entre les deux plans internes est faible et ne peut pas servir de réserve d'énergie pour les appels de courant rapides. Le découplage local décrit à la section~\ref{sec:schema} porte donc seul cette fonction.

\subsection{Placement}

Le placement a été conduit du plus contraignant au plus libre. Les composants volumineux et ceux accessibles depuis l'extérieur du boîtier ont été positionnés en premier, puisque leur position est dictée par l'usage de l'appareil et non par des considérations électriques, et le reste du placement s'est ensuite organisé autour d'eux. Un croquis du produit, reproduit à la figure~\ref{fig:croquis_produit}, a servi de point de départ à cette étape en fixant l'encombrement extérieur ainsi que la position de la fenêtre d'écran et des commandes avant tout travail sur la carte.

\begin{wrapfigure}{l}{0.45\linewidth}
    \centering
    \includegraphics[width=1\linewidth]{\assetsdir/croquis_produit.png}
    \caption{Croquis de la face avant du produit, fixant les dimensions extérieures du PCB et la position de l'écran et des commandes}
    \label{fig:croquis_produit}
\end{wrapfigure}

La première contrainte issue de ce croquis est un partage des faces. L'écran et les boutons doivent se trouver du côté de l'utilisateur, ce qui les place sur le Bottom Layer de la carte. Cette face du circuit imprimé correspond donc à la face avant du produit, et le Top Layer à sa face arrière, convention conservée dans tout ce qui suit afin que le texte se lise directement sur les vues de l'annexe~\ref{ann:layout}.

Le Bottom Layer accueille ainsi l'écran, qui en occupe la plus grande partie, ses trous de fixation, les six boutons de commande et leurs résistances de rappel. Le boîtier entourant l'ensemble, cette face n'est pas pour autant interdite aux autres composants, et quelques circuits y sont reportés là où la place manque sur l'autre face. Le Top Layer, libre de toute contrainte ergonomique, reçoit la quasi-totalité de l'électronique et peut donc être organisé selon des critères purement électriques.

La position des éléments traversant une paroi du boîtier a été décidée de la même manière. La prise casque occupe l'angle inférieur gauche et le connecteur USB-C le bord inférieur droit du Top Layer. Le lecteur de carte microSD et le potentiomètre de volume sont répartis sur le bord gauche, égamelement sur le Top Layer, où ils restent atteignables sans traverser la zone d'écran.

Le Top Layer se divise alors en trois zones. La partie inférieure gauche regroupe la chaîne analogique, du convertisseur, au filtre, à l'amplificateur puis à la prise casque, ce qui maintient les liaisons différentielles courtes et les garde à l'écart du reste de la carte. La bande centrale et droite rassemble l'alimention, de l'entrée USB-C au chargeur, à la jauge, au connecteur de batterie et aux étages à découpage. La partie supérieure porte le numérique, avec le module micro-contrôleur, le lecteur de carte SD et l'expanseur d'entrées-sorties.

Cette répartition ne suffit pas à isoler l'analogique du reste. La surface de la carte étant fixée par le boîtier et la position des connecteurs par la mécanique, les étages à découpage restent voisins de la zone analogique et aucun placement ne permettait de les en écarter davantage. La distance ne suffit donc pas à elle seule. La zone analogique reçoit pour cette raison un plan de masse propre sur le Top Layer, relié au plan interne par des vias répartis sur toute sa surface, qui offre à chaque courant de retour un chemin situé immédiatement à côté de la piste aller au lieu de le contraindre à descendre vers le plan interne puis à remonter. La boucle formée par l'aller et le retour s'en trouve réduite, ce qui diminue à la fois ce que ces liaisons rayonnent et ce qu'elles captent~\cite{prof_bogatin_rules_2013}.

À l'intérieur de chaque zone, les réseaux résistifs de configuration du chargeur et du convertisseur abaisseur sont regroupés au plus près de leurs circuits, ces liaisons peu chargées en courant étant sensibles au bruit capté, qui s'y traduirait directement par une erreur de configuration.

Le placement réserve enfin une trentaine de points de test répartis sur les deux faces, qui rendent les alimentations, les bus et les principales lignes de commande accessibles à la sonde sans démontage.